\section{Planteamiento del problema}

La simulación de sistemas cuánticos en computadoras clásicas está condicionada por el crecimiento exponencial de la memoria requerida. Representar un estado cuántico de $n$ qubits implica almacenar $2^{n}$ números complejos, lo que rápidamente supera la capacidad de los sistemas de memoria convencionales. En este contexto, técnicas como la poda de estados y la compresión con pérdida reducen el consumo de recursos, pero lo hacen modificando la representación numérica del estado y, por tanto, alterando el régimen de exactitud de la simulación.

La simulación exacta sigue siendo necesaria para validar algoritmos cuánticos, contrastar implementaciones y estudiar su comportamiento antes de ejecutarlos en hardware cuántico. Por ello, el problema no consiste solo en reducir memoria, sino en hacerlo sin introducir discrepancias en las amplitudes del estado cuando la exactitud completa es un requisito. Bajo esta formulación, resulta pertinente estudiar estrategias de compresión sin pérdida que permitan mejorar el uso de recursos computacionales preservando la igualdad exacta frente a una ejecución de referencia no comprimida.

Con base en este problema, se plantean los objetivos que orientan el desarrollo del trabajo.
